\subsection{张量积与同伦}

命题3. --- 设C, C \(_{1}\) 为右A-模的两个复形，C'，C' \(_{1}'\) 为左A-模的两个复形，以及u：C → C \(_{1}\) ，v：C → C \(_{1}\) ，u' : C' → C \(_{1}'\) ，v' : C' → C \(_{1}'\) 是复形的态射。

\begin{enumerate}
\def\labelenumi{\alph{enumi})}
\item
  若 \(u\) 与 \(u'\) 分别同伦于 \(v\) 与 \(v'\)，则两个态射 \(u \otimes u'\) 与 \(v \otimes v'\)（从 \(\mathbf{C} \otimes_{\mathbf{A}} \mathbf{C}'\) 到 \(\mathbf{C}_1 \otimes_{\mathbf{A}} \mathbf{C}_1'\)）是同伦的。
\item
  如果 u 和 \(u'\) 是同伦等价， \(u \otimes u'\) 是一个同伦等价。
\item
  如果 C 或 \(\mathbf{C}'\) 同伦于零，则 \(\mathbf{C} \otimes_{\mathbf{A}} \mathbf{C}'\) 同伦于零。
\end{enumerate}

让我们用同一个字母 d 表示复形 C, \(C_{1}\), \(C'\), \(C_{1}'\) 的微分，并用 D 表示复形 \(C \otimes_{A} C'\) 与 \(C_{1} \otimes_{A} C_{1}'\) 的微分。

若 \(u\) (相应地， \(u'\) ) 同伦于 \(v\) (相应地， \(v'\) )，则存在一个次数为 1 的分次同态 \(s: \mathbf{C} \to \mathbf{C}_1\) (相应地， \(s': \mathbf{C}' \to \mathbf{C}_1'\) ）使得

\[
u - v = d s + s d \qquad (\mathrm{resp.} u ^ {\prime} - v ^ {\prime} = d s ^ {\prime} + s ^ {\prime} d).\tag{2}
\]

设 \(\mathbf{S}:\mathbf{C}\otimes_{\mathbf{A}}\mathbf{C}^{\prime}\to \mathbf{C}_{1}\otimes_{\mathbf{A}}\mathbf{C}_{1}^{\prime}\) 是唯一的次数为1的分次同态，使得对于 \(x\in \mathbf{C}_p\) , \(y\in \mathbf{C}_q^\prime\) , \(p,q\in \mathbf{Z}\) ，有

\[
\mathbf {S} (x \otimes y) = s (x) \otimes u ^ {\prime} (y) + (- 1) ^ {p} v (x) \otimes s ^ {\prime} (y).\tag{3}
\]

于是我们有，采用前述记号：

\[
\begin{array}{r l} (\mathrm{DS} + \mathrm{SD}) (x \otimes y) & = \mathrm{D} (s x \otimes u ^ {\prime} y) + (- 1) ^ {p} \mathrm{D} (v x \otimes s ^ {\prime} y) + \mathrm{S} (d x \otimes y) + \\ & + (- 1) ^ {p} \mathrm{S} (x \otimes d y) = \\ & = d s x \otimes u ^ {\prime} y + (- 1) ^ {p + 1} s x \otimes d u ^ {\prime} y + (- 1) ^ {p} d v x \otimes s ^ {\prime} y + v x \otimes d s ^ {\prime} y + \\ & + s d x \otimes u ^ {\prime} y + (- 1) ^ {p - 1} v d x \otimes s ^ {\prime} y + (- 1) ^ {p} s x \otimes u ^ {\prime} d y + v x \otimes s ^ {\prime} d y \\ & = (d s + s d) (x) \otimes u ^ {\prime} y + v x \otimes (d s ^ {\prime} + s ^ {\prime} d) (y) \\ & = (u x - v x) \otimes u ^ {\prime} y + v x \otimes (u ^ {\prime} y - v ^ {\prime} y) = u x \otimes u ^ {\prime} y - v x \otimes v ^ {\prime} y. \end{array}
\]

这给出 DS + SD = u ⊗ u' - v ⊗ v'，由此得 a)。

让我们证明b)。如果u和 \(u'\) 是同伦等价，则存在复形的同态 \(\alpha: C_{1} \to C\) 与 \(\alpha': C_{1}' \to C'\) 使得 \(u \circ \alpha\) , \(\alpha \circ u\) , \(u' \circ \alpha'\) , \(\alpha' \circ u'\) 分别同伦于 \(Id_{C_{1}}\) , \(Id_{C}\) , \(Id_{C_{1}}\) 与 \(Id_{C'}\) 。而 \((u \otimes u') \circ (\alpha \otimes \alpha')\) 等于 \((u \circ \alpha) \otimes (u' \circ \alpha')\) ，由a)它同伦于 \(Id_{C_{1}} \otimes Id_{C_{1}} = Id_{C_{1} \otimes C_{1}}\) ；而 \((\alpha \otimes \alpha') \circ (u \otimes u')\) 同伦于 \(Id_{C \otimes C'}\) ，由此得 b)。最后，c) 由 b) 应用于 \(C_{1}\) 或 \(C_{1}'\) 为零的情形得出。

推论 1. --- 设 C' 是左 A-模的一个分裂复形，使得 H(C') 是平坦的。对于右 A-模的每个复形 C，典范映射

\[
\gamma (\mathrm{C}, \mathrm{C} ^ {\prime}): \mathrm{H} (\mathrm{C}) \otimes_ {\mathrm{A}} \mathrm{H} (\mathrm{C} ^ {\prime}) \rightarrow \mathrm{H} (\mathrm{C} \otimes_ {\mathrm{A}} \mathrm{C} ^ {\prime})
\]

是双射。

根据X，第35页，定义6，存在一个同伦等价 \(u': C' \to H(C')\) 。由命题3， \(1_{C} \otimes u': C \otimes_{A} C' \to C \otimes_{A} H(C')\) 是一个同伦等价；由于

\[
\mathrm{H} \left(1 _ {\mathrm{C}} \otimes u ^ {\prime}\right) \circ \gamma (\mathrm{C}, \mathrm{C} ^ {\prime}) = \gamma (\mathrm{C}, \mathrm{H} (\mathrm{C} ^ {\prime})) \circ \left(1 _ {\mathrm{C}} \otimes \mathrm{H} (u ^ {\prime})\right),
\]

由于 \(\mathbf{H}(1_{\mathbf{C}} \otimes u')\) 与 \(\mathbf{H}(u')\) 都是双射的，只需证明 \(\gamma(\mathbf{C}, \mathbf{H}(\mathbf{C}'))\) 是双射的，于是我们归结为 \(\mathbf{C}'\) 是平坦且微分为零的情形。在此情形下，典范正合序列

\begin{enumerate}
\def\labelenumi{(\Roman{enumi})}
\tightlist
\item
\end{enumerate}

\[
0 \to \mathbf {Z} (\mathbf {C}) \xrightarrow {i} \mathbf {C} \xrightarrow {\delta} \mathbf {B} (\mathbf {C}) \to 0\tag{II}
\]

\[
0 \rightarrow \mathrm{B} (\mathrm{C}) \xrightarrow {j} \mathrm{Z} (\mathrm{C}) \xrightarrow {\pi} \mathrm{H} (\mathrm{C}) \rightarrow 0
\]

给出正合序列：

\[
\begin{array}{l} 0 \longrightarrow Z (C) \otimes_ {A} C ^ {\prime} \xrightarrow {i \otimes 1} C \otimes_ {A} C ^ {\prime} \quad \xrightarrow {\delta \otimes 1} B (C) \otimes_ {A} C ^ {\prime} \longrightarrow 0 \\ 0 \longrightarrow B (C) \otimes_ {A} C ^ {\prime} \xrightarrow {j \otimes 1} Z (C) \otimes_ {A} C ^ {\prime} \xrightarrow {\pi \otimes 1} H (C) \otimes_ {A} C ^ {\prime} \longrightarrow 0. \end{array}
\]

由于 \(d = i \circ j \circ \delta\)，我们有 \(\mathbf{D} = d \otimes 1_{\mathbf{C}'} = (i \otimes 1) \circ (j \otimes 1) \circ (\delta \otimes 1)\)，这表明典范映射 \(\mathbf{Z}(\mathbf{C}) \otimes_{\mathbf{A}} \mathbf{C}' \to \mathbf{Z}(\mathbf{C} \otimes_{\mathbf{A}} \mathbf{C}')\) 与 \(\mathbf{B}(\mathbf{C}) \otimes_{\mathbf{A}} \mathbf{C}' \to \mathbf{B}(\mathbf{C} \otimes_{\mathbf{A}} \mathbf{C}')\) 是双射，因此通过取商，\(\gamma(\mathbf{C}, \mathbf{C}')\) 也是双射的。

推论2. --- 设N为平坦的左A-模。对于右A-模的每个复形C，典范同态

\[
\gamma_ {n} (\mathbf {C}, \mathbf {N}): \mathrm{H} _ {n} (\mathbf {C}) \otimes_ {\mathbf {A}} \mathbf {N} \rightarrow \mathrm{H} _ {n} (\mathbf {C} \otimes_ {\mathbf {A}} \mathbf {N})
\]

都是双射的。

推论3. --- 设C'为左A-模的复形，使得B(C')和H(C')是投射的。对于右A-模的每个复形C，映射γ(C, C')是双射。

事实上，C'是分裂的（X，第35页，例4），且H(C')是投射的；因此可以应用推论1。

注记. --- 1) 利用交换同构，由推论1、2和3可推出将张量积的两个参数角色互换后所得的类似结论。

\begin{enumerate}
\def\labelenumi{\arabic{enumi})}
\setcounter{enumi}{1}
\tightlist
\item
  我们将在下文（X，第79页，推论4）看到，当假设C'和H(C')是平坦的且C'在右侧有界时，推论1的结论也成立。
\end{enumerate}

